% Motor: XeLaTeX, fijado en .latexmkrc y en la receta del editor.
\documentclass[aspectratio=169,10pt,c]{beamer}

% !TEX root = slides.tex
% Preambulo comun: estilo UTP, colores, fuentes, plantillas de beamer y
% comandos de apoyo didactico. Se carga desde slides.tex tras \documentclass.

% La opción "c" de beamer reparte el espacio libre 1:1.5 (arriba:abajo), por lo
% que el contenido queda por encima del centro. Igualamos ambos rellenos.
\makeatletter
\define@key{beamerframe}{c}[true]{%
  \beamer@frametopskip=0pt plus 1fill\relax%
  \beamer@framebottomskip=0pt plus 1fill\relax%
  \beamer@frametopskipautobreak=0pt plus .5\paperheight\relax%
  \beamer@framebottomskipautobreak=0pt plus .5\paperheight\relax%
}
\makeatother

\usepackage{graphicx}
\usepackage{tikz}
\usepackage{xcolor}
\usepackage{iftex}

\ifPDFTeX
  \usepackage[utf8]{inputenc}
  \usepackage[T1]{fontenc}
  \usepackage{lmodern}          % tipografia escalable para el respaldo con pdflatex
  \renewcommand{\familydefault}{\sfdefault}
\else
  \usepackage{fontspec}
  \IfFontExistsTF{Arial}
    {\setsansfont{Arial}[BoldFont={Arial Bold}]}
    {\setsansfont{Helvetica}[BoldFont={Helvetica Bold}]}
\fi

\usepackage[spanish,es-nodecimaldot]{babel}
\usepackage{microtype}
\usepackage{booktabs}
\usepackage{amssymb}
\usepackage{array}
\usetikzlibrary{arrows.meta,positioning,fit,calc,shapes.geometric}

\graphicspath{{assets/}}

% ---------- Identidad visual UTP ----------
\definecolor{utpRojo}{HTML}{C42B35}
\definecolor{utpNegro}{HTML}{111111}
\definecolor{utpGris}{HTML}{5F5F5F}
\definecolor{utpGrisMedio}{HTML}{9A9A9A}
\definecolor{utpGrisClaro}{HTML}{F1F1F1}
\definecolor{utpNaranja}{HTML}{C1670B}

\usetheme{default}
\setbeamercolor{normal text}{fg=utpNegro,bg=white}
\setbeamercolor{structure}{fg=utpRojo}
\setbeamercolor{frametitle}{fg=utpNegro,bg=white}
\setbeamercolor{title}{fg=utpNegro}
\setbeamercolor{subtitle}{fg=utpNegro}
\setbeamercolor{author}{fg=utpNegro}
\setbeamercolor{date}{fg=utpNegro}
\setbeamercolor{block title}{fg=white,bg=utpNegro}
\setbeamercolor{block body}{fg=utpNegro,bg=utpGrisClaro}
\setbeamercolor{alertblock title}{fg=white,bg=utpRojo}
\setbeamercolor{alertblock body}{fg=utpNegro,bg=utpRojo!7}
\setbeamercolor{exampleblock title}{fg=white,bg=utpGris}
\setbeamercolor{exampleblock body}{fg=utpNegro,bg=utpGris!8}

\setbeamerfont{title}{series=\bfseries,size=\fontsize{19.5}{22.5}\selectfont}
\setbeamerfont{subtitle}{size=\fontsize{11.5}{13.5}\selectfont}
\setbeamerfont{date}{series=\bfseries,size=\fontsize{11.5}{13.5}\selectfont}
\setbeamerfont{author}{size=\fontsize{10.5}{12.5}\selectfont}
\setbeamerfont{frametitle}{series=\bfseries,size=\large}
\setbeamerfont{block title}{series=\bfseries,size=\small}
\setbeamerfont{block body}{size=\normalsize}

\setbeamertemplate{navigation symbols}{}
\setbeamertemplate{headline}{}
\setbeamertemplate{blocks}[default]
\setbeamertemplate{itemize item}{\color{utpRojo}\small\raise.15ex\hbox{$\blacktriangleright$}}
\setbeamertemplate{itemize subitem}{\color{utpGris}\scriptsize$\bullet$}
\setbeamertemplate{enumerate item}{\color{utpRojo}\insertenumlabel.}
\setbeamertemplate{background canvas}{\color{white}\rule{\paperwidth}{\paperheight}}

% Título de diapositiva: negrita negra a la izquierda, logo UTP a la derecha
\setbeamertemplate{frametitle}{%
  \begin{tikzpicture}[remember picture,overlay]
    \node[anchor=north east,inner sep=0pt]
      at ([xshift=-7mm,yshift=-5.5mm]current page.north east)
      {\includegraphics[width=17mm]{utp-logo.png}};
  \end{tikzpicture}%
  \vspace{1.2mm}%
  \begin{minipage}[t]{.80\textwidth}\raggedright\insertframetitle\end{minipage}\par
  \vspace{1.1mm}{\color{utpRojo}\rule{\textwidth}{1.1pt}}
}

% Solo numeración discreta; la plantilla UTP no lleva pie de página
\setbeamertemplate{footline}{%
  \hfill{\color{utpGrisMedio}\scriptsize\insertframenumber}\hspace{6mm}\vspace{2.5mm}%
}

% Portada UTP
\setbeamertemplate{title page}{%
  \begin{tikzpicture}[remember picture,overlay]
    \node[anchor=center,inner sep=0pt] at (current page.center) {%
      \includegraphics[width=\paperwidth,height=\paperheight]{utp-cover-bg.jpeg}%
    };
    \node[anchor=north,align=center,text=utpNegro,inner sep=0pt,
      text width=145mm,font=\usebeamerfont{title}]
      at ([yshift=-25.2mm]current page.north) {\inserttitle};
    \node[anchor=north,align=center,text=utpNegro,inner sep=0pt,
      text width=142mm,font=\usebeamerfont{subtitle}]
      at ([yshift=-38.2mm]current page.north) {\insertsubtitle};
    \node[anchor=north,align=center,text=utpNegro,inner sep=0pt,
      font=\usebeamerfont{date}]
      at ([yshift=-50.4mm]current page.north) {\insertdate};
    \node[anchor=north,align=center,text=utpNegro,inner sep=0pt,
      font=\usebeamerfont{author}]
      at ([yshift=-58.4mm]current page.north) {\insertauthor};
  \end{tikzpicture}%
}

% ---------- Comandos de apoyo didáctico ----------
\newcommand{\concepto}[2]{%
  \begin{block}{#1}\large #2\end{block}}
\newcommand{\pregunta}[1]{%
  \begin{alertblock}{Decide y explica}#1\end{alertblock}}
\newcommand{\analogia}[1]{%
  \begin{exampleblock}{En palabras simples}\small #1\end{exampleblock}}
\newcommand{\enlarealidad}[1]{%
  \begin{block}{Cómo se ve en la realidad}\small #1\end{block}}
\newcommand{\ojo}[1]{%
  \begin{alertblock}{Ojo con esto}\small #1\end{alertblock}}
\newcommand{\apoyo}{{\normalfont\footnotesize\color{utpGris}~· material de apoyo}}
\newcommand{\marca}[1]{\tikz[baseline=(m.base)]{%
  \node[rounded corners=1.5pt,fill=utpGrisClaro,draw=utpGrisMedio,inner xsep=3pt,inner ysep=1.5pt,
        font=\footnotesize] (m) {#1};}}

\tikzset{
  caja/.style={draw=utpNegro, very thick, rounded corners=1.5pt, fill=white,
    align=center, minimum height=0.72cm, inner sep=5pt, font=\small},
  cajaazul/.style={caja, draw=utpGris, fill=white},
  cajaverde/.style={caja, draw=utpNegro, fill=utpGrisClaro},
  cajanaranja/.style={caja, draw=utpRojo!75, fill=utpRojo!7},
  flecha/.style={-{Stealth[length=2.3mm]}, very thick, draw=utpNegro},
  flechavuelta/.style={{Stealth[length=2.3mm]}-{Stealth[length=2.3mm]}, very thick, draw=utpNegro}
}

% ---------- Notación pata de gallo (crow's foot) ----------
% La rotación indica hacia dónde está la entidad: 0 derecha, 180 izquierda,
% 90 arriba, -90 abajo. La pata se abre hacia la entidad del lado "muchos".
\tikzset{
  relacion/.style={thick, draw=utpNegro},
  pics/uno/.style={code={%
    \draw[thick,utpNegro,line cap=round] (0,-1.5mm) -- (0,1.5mm);}},
  pics/muchos/.style={code={%
    \draw[thick,utpNegro,line cap=round,line join=round]
      (2.6mm,1.5mm) -- (0,0) -- (2.6mm,-1.5mm);
    \draw[thick,utpNegro,line cap=round] (0,0) -- (2.6mm,0);}},
}
\newcommand{\simbuno}{\tikz[baseline=-0.6ex]{\draw[relacion](0,0)--(7mm,0);\pic at (4.4mm,0){uno};}}
\newcommand{\simbmuchos}{\tikz[baseline=-0.6ex]{\draw[relacion](0,0)--(7mm,0);\pic at (4.4mm,0){muchos};}}


\title{Del modelo conceptual al modelo físico}
\subtitle{Orden de SELECT, modelos de datos y\\ Forward Engineering y Reverse Engineering}
\date{Semana 10 - Sesión 20}
\author{Rel Guzman}

% Espacio en blanco para completar
\newcommand{\bl}[1][1.9cm]{\underline{\hspace{#1}}}
% Nombre de tabla, siempre en naranja
\newcommand{\tbl}[1]{\textcolor{utpNaranja}{\texttt{#1}}}
% Palabra reservada de SQL dentro de un bloque de código
\newcommand{\kw}[1]{\textcolor{utpRojo}{\textbf{#1}}}
% Parte nueva de un bloque de codigo que crece paso a paso
\newcommand{\nuevo}[1]{{%
  \setlength{\fboxsep}{1.2pt}%
  \begingroup
  \renewcommand{\kw}[1]{\textcolor{utpNegro}{\textbf{##1}}}%
  \colorbox{yellow!45}{\textcolor{utpNegro}{#1}}%
  \endgroup}}
% Comentario dentro de un bloque de código SQL
\newcommand{\com}[1]{\textcolor{utpGris}{-- #1}}
% Marcas de clave primaria y clave foránea
\newcommand{\pk}{{\scriptsize\color{utpRojo}\textbf{PK}}}
\newcommand{\fk}{{\scriptsize\color{utpNaranja}\textbf{FK}}}
% Insignia de práctica
\newcommand{\ins}[1]{\tikz[baseline=-0.5ex]{\node[rounded corners=2pt,fill=utpRojo,text=white,
  font=\bfseries\scriptsize,inner xsep=5pt,inner ysep=2.5pt] {#1};}}
% Etiqueta oscura: nombre del gestor o de la forma de escribir algo
\newcommand{\etiqueta}[1]{\tikz[baseline=-0.5ex]{\node[rounded corners=2pt,fill=utpNegro,text=white,
  font=\bfseries\scriptsize,inner xsep=6pt,inner ysep=3pt] {#1};}}
% Celda resaltada en rojo claro
\newcommand{\mal}[1]{{\setlength{\fboxsep}{1.2pt}\colorbox{utpRojo!18}{#1}}}
% Barra invertida de las rutas de Windows
\newcommand{\bs}{\textbackslash{}}
% Bloque de código sobre fondo gris
\newsavebox{\sqlbuf}
\newenvironment{sqlcode}[1][.86\textwidth]
  {\begin{center}\begin{lrbox}{\sqlbuf}\begin{minipage}{#1}\ttfamily\footnotesize\vspace{6pt}\raggedright}
  {\vspace{6pt}\end{minipage}\end{lrbox}\colorbox{utpGrisClaro}{\usebox{\sqlbuf}}\end{center}}
% Sangría dentro de un bloque de código
\newcommand{\ind}{\hspace*{1.6em}}
% Anatomía de una instrucción: cada pieza del código con su pregunta al lado
\tikzset{
  pieza/.style={fill=utpGrisClaro,font=\ttfamily\footnotesize,anchor=west,inner sep=5pt,
    minimum height=0.62cm},
  para/.style={font=\small,anchor=west,text=utpNegro}
}
\newcommand{\piezax}{7.6}
\newcommand{\pieza}[3]{%
  \node[pieza] at (0,#1) {#2};
  \node[para] at (\piezax,#1) {\textcolor{utpRojo}{$\blacktriangleleft$}\ #3};}
% Diapositiva de sección: una palabra bien grande (\seccion[0]{...} la deja sin línea roja)
\newcommand{\seccion}[2][1]{%
\begin{frame}[plain]
  \begin{tikzpicture}[remember picture,overlay]
    \node[anchor=north east,inner sep=0pt]
      at ([xshift=-10mm,yshift=-8mm]current page.north east)
      {\includegraphics[width=43mm]{utp-logo.png}};
    \node[anchor=center,inner sep=0pt] at (current page.center)
      {\resizebox{100mm}{!}{\bfseries\color{utpNegro}#2}};
    \ifnum#1=1
    \node[anchor=north,inner sep=0pt] at ([yshift=-11mm]current page.center)
      {\color{utpRojo}\rule{100mm}{1.6pt}};
    \fi
  \end{tikzpicture}
\end{frame}}
% Cabecera de columna dentro de una tabla en monoespaciado
\newcommand{\cab}[1]{\normalfont\bfseries #1}
% Resultado de una consulta: \res{columnas}{cabecera}{filas}
\newcommand{\res}[3]{%
  \centering{\scriptsize\color{utpGris}resultado}\\[3pt]
  \footnotesize\ttfamily
  \begin{tabular}{#1}
  \toprule
  #2 \\
  \midrule
  #3
  \bottomrule
  \end{tabular}}
% Columna del modelo lógico: nombre y, en gris, su tipo
\newcommand{\col}[2]{\texttt{#1}~{\color{utpGris}#2}}
\begin{document}
{\setbeamertemplate{footline}{}\maketitle}
\begin{frame}[c]
\begin{block}{Logro de la sesión}
Al finalizar, podrás explicar el \textbf{orden lógico de un SELECT}, distinguir los modelos \textbf{conceptual, lógico y físico}, y reconocer cómo \textbf{Forward Engineering} y \textbf{Reverse Engineering} relacionan el diseño con la base de datos. También compararás el modelo físico en \textbf{SQL Server, MySQL y SQLite}.
\end{block}
\end{frame}
\begin{frame}{Revisión de la sesión anterior: orden de SELECT}
\begin{center}
\begin{tikzpicture}[y=0.54cm]
  \renewcommand{\piezax}{7.4}
  \pieza{7}{\kw{SELECT} c.nombre, \kw{COUNT}(*) \kw{AS} pedidos}{qué columnas o cálculos se \textbf{muestran}}
  \pieza{6}{\kw{FROM} cliente c}{la tabla \textbf{principal}}
  \pieza{5}{\kw{JOIN} pedido p \kw{ON} c.id = p.cliente\_id}{qué otra tabla se \textbf{une}, y por dónde}
  \pieza{4}{\kw{WHERE} p.fecha >= '2026-01-01'}{qué \textbf{filas} entran}
  \pieza{3}{\kw{GROUP BY} c.nombre}{cómo se \textbf{agrupan}}
  \pieza{2}{\kw{HAVING} \kw{COUNT}(*) >= 2}{qué \textbf{grupos} se quedan}
  \pieza{1}{\kw{ORDER BY} pedidos \kw{DESC}}{en qué \textbf{orden} salen}
  \pieza{0}{\kw{LIMIT} 5;}{cuántas filas, \textbf{como máximo}}
\end{tikzpicture}
\end{center}
\vspace{-0.05cm}
\begin{itemize}\setlength\itemsep{3pt}
  \item \texttt{SELECT} indica qué mostrar. \texttt{FROM} se usa al consultar tablas. Lo demás se agrega si hace falta.
  \item Si se usan, se escriben \textbf{siempre en este orden}.
  \item \texttt{SELECT DISTINCT} quita las filas \textbf{repetidas} del resultado.
\end{itemize}
\vspace{0.05cm}
\small\color{utpGris} Se lee: «los 5 clientes con más pedidos este año, entre los que tienen 2 o más».
\end{frame}
\begin{frame}{Se escribe en un orden, se resuelve en otro}
\small Este es el \textbf{orden lógico} de la consulta. El motor puede usar otro plan físico para ejecutarla:
\vspace{0.25cm}
\begin{center}
\resizebox{\textwidth}{!}{%
\begin{tikzpicture}[font=\small]
  \tikzset{paso/.style={caja,text width=1.75cm,minimum height=1.0cm,font=\small\ttfamily\bfseries}}
  \node[paso,cajanaranja] (a) {FROM\\[-1pt]JOIN};
  \node[paso,right=4mm of a] (b) {WHERE};
  \node[paso,right=4mm of b] (c) {GROUP BY};
  \node[paso,right=4mm of c] (d) {HAVING};
  \node[paso,cajanaranja,right=4mm of d] (e) {SELECT};
  \node[paso,right=4mm of e] (f) {ORDER BY};
  \node[paso,right=4mm of f] (g) {LIMIT};
  \foreach \u/\v in {a/b,b/c,c/d,d/e,e/f,f/g} \draw[flecha] (\u)--(\v);
  \foreach \n/\t/\q in {a/1/junta las tablas,b/2/descarta filas,c/3/arma los grupos,
                     d/4/descarta grupos,e/5/elige las columnas,f/6/ordena,g/7/corta} {
    \node[font=\scriptsize\bfseries,text=utpRojo,above=1.5mm of \n] {\t};
    \node[font=\scriptsize,text=utpGris,below=1.5mm of \n,text width=1.8cm,align=center] {\q};
  }
\end{tikzpicture}}
\end{center}
\vspace{0.2cm}
\begin{itemize}\setlength\itemsep{6pt}
  \item Por eso \texttt{WHERE COUNT(*) >= 2} da error: en el paso 2 los grupos \textbf{todavía no existen}.
  \item Y por eso \texttt{ORDER BY pedidos} funciona: en el paso 6 el \texttt{SELECT} \textbf{ya puso} ese nombre.
\end{itemize}
\vspace{0.1cm}
\ojo{Si aparece \texttt{DISTINCT}, quita duplicados después de \texttt{SELECT} y antes de \texttt{ORDER BY}. Aquí usamos \texttt{LIMIT}, como en MySQL y SQLite.}
\end{frame}

% Repaso: alias de agregación y diferencia entre WHERE y HAVING.
\begin{frame}{¿Por qué da error \texttt{WHERE total > 1}?}
\small Queremos mostrar los desarrolladores que tienen \textbf{más de un videojuego}.
\vspace{0.2cm}
\begin{columns}[c,onlytextwidth]
\column{.49\textwidth}
\begin{sqlcode}[\dimexpr\textwidth-2\fboxsep\relax]
\kw{SELECT} developer, \kw{COUNT}(*) \kw{AS} total\\
\kw{FROM} videogames\\
\mal{\kw{WHERE} total > 1}\\
\kw{GROUP BY} developer;
\end{sqlcode}
\column{.47\textwidth}
\small
\begin{itemize}\setlength\itemsep{7pt}
  \item \texttt{AS total} pone un \textbf{nombre al resultado del conteo}. No crea una columna en \texttt{videogames}.
  \item \texttt{WHERE} revisa los videojuegos \textbf{fila por fila}, antes de formar los grupos.
  \item En ese paso todavía no hay un conteo por desarrollador que se pueda comparar con 1.
\end{itemize}
\end{columns}
\vspace{0.2cm}
\ojo{Cambiarlo por \texttt{WHERE COUNT(*) > 1} \textbf{también da error}: un conteo del grupo no se filtra en \texttt{WHERE}.}
\vspace{0.1cm}
\footnotesize\color{utpGris} Esta consulta falla también en SQLite. Aunque permite algunos alias en \texttt{WHERE}, no permite allí este alias de \texttt{COUNT(*)}.
\note{Fuentes: https://dev.mysql.com/doc/refman/8.0/en/problems-with-alias.html ; https://www.postgresql.org/docs/current/sql-select.html ; https://www.sqlite.org/lang_select.html . Consulta comprobada con SQLite, sobre una tabla videogames que solo contiene developer.}
\end{frame}

\begin{frame}{¿Por qué sí funciona \texttt{HAVING COUNT(*) > 1}?}
\begin{columns}[c,onlytextwidth]
\column{.49\textwidth}
\begin{sqlcode}[\dimexpr\textwidth-2\fboxsep\relax]
\kw{SELECT} developer, \kw{COUNT}(*) \kw{AS} total\\
\kw{FROM} videogames\\
\kw{GROUP BY} developer\\
\nuevo{\kw{HAVING COUNT}(*) > 1};
\end{sqlcode}
\column{.47\textwidth}
\small\textbf{Ejemplo: 6 videojuegos en la tabla}
\vspace{0.15cm}

\renewcommand{\arraystretch}{1.3}
\begin{tabular}{@{}lcl@{}}
\toprule
\textbf{Desarrollador} & \textbf{Conteo} & \textbf{¿Se queda?} \\
\midrule
Aurora & 3 & Sí: $3 > 1$ \\
Pixel & 1 & No: $1 \not> 1$ \\
Nube & 2 & Sí: $2 > 1$ \\
\bottomrule
\end{tabular}
\end{columns}
\vspace{0.2cm}
\small
\begin{enumerate}\setlength\itemsep{5pt}
  \item \texttt{GROUP BY developer} reúne los juegos del mismo desarrollador.
  \item \texttt{HAVING COUNT(*) > 1} cuenta los juegos de cada grupo y conserva los grupos con más de uno.
  \item \texttt{SELECT} muestra el desarrollador y su conteo, con el nombre \texttt{total}: \textbf{Aurora, 3} y \textbf{Nube, 2}. El orden no está garantizado sin \texttt{ORDER BY}.
\end{enumerate}
\vspace{0.1cm}
\concepto{La diferencia}{\small \texttt{WHERE} decide \textbf{qué filas entran al grupo}. \texttt{HAVING} decide \textbf{qué grupos aparecen en el resultado}.}
\note{Fuentes: https://www.postgresql.org/docs/current/sql-select.html ; https://www.sqlite.org/lang_select.html . Se escribe COUNT(*) en HAVING para no depender del soporte de alias de cada gestor.}
\end{frame}

\seccion[0]{Modelos de datos}
% ---------------------------------------------------------------------
\begin{frame}{Un mismo diseño, en tres niveles de detalle}
\begin{center}
\resizebox{\textwidth}{!}{%
\begin{tikzpicture}[font=\small]
  \tikzset{nivel/.style={caja,text width=2.3cm,minimum height=1.1cm}}
  \node[nivel,cajaazul] (a) {Caso de\\estudio};
  \node[nivel,right=6mm of a] (b) {\textbf{Modelo}\\\textbf{conceptual}};
  \node[nivel,right=6mm of b] (c) {\textbf{Modelo}\\\textbf{lógico}};
  \node[nivel,cajanaranja,right=6mm of c] (d) {\textbf{Modelo}\\\textbf{físico}};
  \node[nivel,cajaverde,right=6mm of d] (e) {Optimización};
  \foreach \u/\v in {a/b,b/c,c/d,d/e} \draw[flecha] (\u)--(\v);
  \foreach \n/\q/\f in {a/el problema\\en palabras/1\_contexto.md,
                        b/qué cosas hay y\\cómo se relacionan/2\_modelo\_conceptual.dbml,
                        c/{tablas, columnas,\\PK y FK}/3\_modelo\_logico.dbml,
                        d/el SQL exacto\\de un gestor/4\_modelo\_fisico.sql,
                        e/índices para que\\sea rápido/\textrm{\color{utpGris}después}} {
    \node[font=\scriptsize,below=2mm of \n,text width=2.9cm,align=center] (q\n) {\q};
    \node[font=\tiny\ttfamily,text=utpRojo,below=0.5mm of q\n,text width=2.7cm,align=center] {\f};
  }
\end{tikzpicture}}
\end{center}
\vspace{0.1cm}
\begin{itemize}\setlength\itemsep{4pt}
  \item Cada modelo \textbf{agrega detalle} al anterior. No se tira nada: se completa.
  \item Los archivos están en la carpeta \texttt{modelado/} de esta sesión.
\end{itemize}
\vspace{0.1cm}
\analogia{Es construir una casa: primero el \textbf{dibujo} de qué ambientes hay, luego el \textbf{plano} con medidas y, al final, los \textbf{planos de obra} con los materiales reales. Ya viviendo en ella, se hacen \textbf{mejoras}.}
\end{frame}
% ---------------------------------------------------------------------
\begin{frame}{El caso: una peluquería de mascotas}
\begin{columns}[T,onlytextwidth]
\column{.58\textwidth}
\small
\begin{itemize}\setlength\itemsep{6pt}
  \item Se registran \textbf{clientes} y sus \textbf{mascotas}.
  \item Una \textbf{cita} es de una mascota, para un \textbf{servicio}, con un \textbf{empleado}, en una fecha y hora. Tiene un estado.
  \item Cada servicio tiene una \textbf{duración} y un \textbf{precio}.
  \item Al completarse la cita se registra su \textbf{pago}: monto y método.
  \item El \textbf{encargado} inicia sesión con usuario y contraseña.
\end{itemize}
\column{.38\textwidth}
\begin{exampleblock}{Fuera de alcance}
\small Varias sedes, inventario, recordatorios, pagos parciales, varias cuentas o roles.
\end{exampleblock}
\end{columns}
\vspace{0.35cm}
\concepto{El primer paso}{\small Los \textbf{sustantivos importantes} del problema (en negrita) son candidatos a \textbf{entidades}. Lo que se dice de ellos, a \textbf{atributos}. Los verbos que los unen, a \textbf{relaciones}.}
\end{frame}
% ---------------------------------------------------------------------
\begin{frame}{Modelo conceptual: qué hay y cómo se relaciona}
\small Describe el negocio con \textbf{entidades}, \textbf{atributos} y \textbf{relaciones}.
\vspace{0.2cm}
\begin{itemize}\setlength\itemsep{9pt}
  \item \textbf{Entidades:} cliente, mascota, cita, servicio, empleado, pago y encargado.
  \item \textbf{Atributos:} el cliente tiene nombre y teléfono. La mascota tiene nombre, especie y raza. La cita tiene fecha, hora y estado.
  \item \textbf{Relaciones:} un cliente es dueño de muchas mascotas. Una mascota puede tener muchas citas. Cada cita corresponde a un servicio y un empleado.
  \item \textbf{Regla del negocio:} al completar una cita se registra su pago total, con monto y método.
\end{itemize}
\vspace{0.2cm}
\concepto{Qué define este nivel}{\small Explica \textbf{qué información necesita la peluquería} y cómo se relaciona. Todavía no elige tipos SQL, claves foráneas ni gestor.}
\end{frame}
% ---------------------------------------------------------------------
\begin{frame}{Modelo lógico: tablas, columnas y claves}
\small Organiza la información en \textbf{tablas} y define cómo identificar y relacionar sus filas.
\vspace{0.2cm}
\begin{itemize}\setlength\itemsep{9pt}
  \item \textbf{Tablas y columnas:} \tbl{cliente} guarda \texttt{id}, \texttt{nombre} y \texttt{telefono}. \tbl{mascota} agrega \texttt{cliente\_id} a sus datos.
  \item \textbf{Clave primaria (PK):} \texttt{id} identifica cada fila de una tabla.
  \item \textbf{Clave foránea (FK):} \texttt{mascota.cliente\_id} referencia a \texttt{cliente.id}. La tabla \tbl{cita} referencia a mascota, servicio y empleado.
  \item \textbf{Dominios y reglas:} se distingue texto, números y fechas. Se indica qué datos son obligatorios y cuáles no se pueden repetir.
\end{itemize}
\vspace{0.2cm}
\concepto{Qué define este nivel}{\small Establece la \textbf{estructura y las restricciones}. En el físico se eligen los tipos y la sintaxis exacta de cada gestor.}
\end{frame}
% ---------------------------------------------------------------------
\begin{frame}{Conceptual, lógico y físico: lado a lado}
\begin{center}
\small
\renewcommand{\arraystretch}{1.25}
\begin{tabular}{@{}>{\raggedright\bfseries}p{2.1cm}>{\raggedright\arraybackslash}p{4.3cm}>{\raggedright\arraybackslash}p{4.3cm}>{\centering\arraybackslash}p{1.8cm}@{}}
\toprule
& \bfseries Qué muestra & \bfseries La cita, en ese modelo & \bfseries ¿Depende del gestor? \\
\midrule
Conceptual & entidades, atributos y relaciones con verbos & «una mascota \textit{tiene} muchas citas» & No \\
Lógico & tablas y columnas;\newline claves, tipos y reglas & \texttt{cita(id \pk, fecha\_hora, \ldots, mascota\_id \fk)} & No \\
Físico & el SQL exacto que crea la base & \texttt{CREATE TABLE cita (\ldots);} & \textbf{Sí} \\
\bottomrule
\end{tabular}
\end{center}
\vspace{0.3cm}
\concepto{La idea clave}{\small El conceptual y el lógico son \textbf{los mismos} para SQL Server, MySQL o SQLite. Solo el \textbf{físico} cambia según el gestor.}
\vspace{0.1cm}
\small En el curso: conceptual y lógico en \textbf{DBML} (dbdiagram), físico en un archivo \texttt{.sql}.
\end{frame}
% =======================================================================
\seccion{Modelo físico}
\begin{frame}{Modelo físico: el SQL que crea la base de datos}
\begin{columns}[c,onlytextwidth]
\column{.58\textwidth}
\begin{sqlcode}[\dimexpr\textwidth-2\fboxsep\relax]
\scriptsize
\kw{CREATE TABLE} cita (\\
\ind id \nuevo{INTEGER \kw{PRIMARY KEY}},\\
\ind fecha\_hora TEXT \kw{NOT NULL},\\
\ind estado TEXT \kw{NOT NULL} \nuevo{\kw{DEFAULT} 'programada'}\\
\ind\ind \nuevo{\kw{CHECK} (estado \kw{IN} ('programada', 'completada',}\\
\ind\ind \nuevo{\hspace*{4.6em}'cancelada'))},\\
\ind mascota\_id \ INTEGER \kw{NOT NULL},\\
\ind servicio\_id INTEGER \kw{NOT NULL},\\
\ind empleado\_id INTEGER \kw{NOT NULL},\\
\ind \nuevo{\kw{FOREIGN KEY} (mascota\_id) \kw{REFERENCES} mascota(id)},\\
\ind \com{\ldots{} y lo mismo para servicio y empleado}\\
\ind \com{un empleado no atiende dos citas a la misma hora}\\
\ind \nuevo{\kw{UNIQUE} (empleado\_id, fecha\_hora)}\\
);
\end{sqlcode}
\column{.38\textwidth}
\small
\begin{itemize}\setlength\itemsep{7pt}
  \item Es el archivo \texttt{4\_modelo\_fisico.sql}.
  \item Está escrito para \textbf{un gestor}: aquí, SQLite.
  \item Aquí las reglas se vuelven \textbf{de verdad}: la base de datos \textbf{rechaza} lo que no las cumple.
\end{itemize}
\vspace{0.2cm}
\enlarealidad{Con este archivo, cualquiera crea la misma base de datos \textbf{vacía} en su computadora.}
\end{columns}
\end{frame}

% Fuentes: https://dev.mysql.com/doc/workbench/en/wb-forward-engineering.html
\begin{frame}{Forward Engineering: del modelo a la base de datos}
\concepto{Ingeniería directa}{\small Parte de un \textbf{modelo definido para un gestor} y genera el SQL que crea su estructura. Puede exportar un script o ejecutarlo en un servidor.}
\vspace{0.2cm}
\begin{enumerate}\setlength\itemsep{5pt}
  \item Definir tablas, tipos, claves y restricciones en el modelo.
  \item Generar y revisar las instrucciones \texttt{CREATE TABLE} y los demás objetos.
  \item Ejecutar el script para crear la estructura de la base de datos.
\end{enumerate}
\vspace{0.2cm}
\enlarealidad{En MySQL Workbench, el modelo de la peluquería permite generar un script para MySQL. Al ejecutarlo se crean tablas como \texttt{cliente}, \texttt{mascota} y \texttt{cita}.}
\vspace{0.15cm}
\small Generar las tablas \textbf{no llena automáticamente los datos} de clientes o citas.
\end{frame}
% Fuentes: https://dev.mysql.com/doc/workbench/en/wb-reverse-engineering.html
\begin{frame}{Reverse Engineering: de la base de datos al modelo}
\concepto{Ingeniería inversa}{\small Parte de una \textbf{base de datos existente} o de su script de creación y reconstruye un modelo de su estructura.}
\vspace{0.2cm}
\begin{enumerate}\setlength\itemsep{5pt}
  \item Leer las tablas, columnas, claves y restricciones existentes.
  \item Importar esa estructura en una herramienta de modelado.
  \item Revisar el modelo obtenido para documentar o modificar la base.
\end{enumerate}
\vspace{0.2cm}
\enlarealidad{Si recibimos la base MySQL de la peluquería sin su modelo, Workbench puede importar el esquema y mostrar las relaciones declaradas con claves foráneas.}
\vspace{0.15cm}
\ojo{Recupera lo que está definido en la base. Las reglas que solo conoce la aplicación y el significado del negocio requieren revisión humana.}
\end{frame}
\begin{frame}{MySQL, SQL Server y SQLite: resumen}
\small
\begin{description}\setlength\itemsep{16pt}
  \item[MySQL] Gestor relacional que funciona como \textbf{servidor}. \textbf{InnoDB} es su motor de almacenamiento predeterminado: admite transacciones y claves foráneas. \textbf{MySQL Workbench} permite administrar bases y trabajar con modelos.
  \item[SQL Server] Gestor relacional de \textbf{Microsoft}, también con servidor. Usa \textbf{Transact-SQL (T-SQL)}. \textbf{SSMS} es una herramienta para conectarse, ejecutar consultas y administrar las bases.
  \item[SQLite] Motor integrado en la propia aplicación, \textbf{sin un servidor separado}. Guarda la base en un archivo local. \textbf{DB Browser for SQLite} permite abrirla y trabajar con sus tablas.
\end{description}
\vspace{0.3cm}
\analogia{El gestor ejecuta SQL y administra los datos. Workbench, SSMS y DB Browser son programas que permiten trabajar con esos gestores.}
\end{frame}
\begin{frame}{Los tres gestores, lado a lado}
\begin{center}
\small
\renewcommand{\arraystretch}{1.2}
\begin{tabular}{@{}>{\raggedright\bfseries}p{2.9cm}>{\raggedright\arraybackslash}p{3.2cm}>{\raggedright\arraybackslash}p{3.2cm}>{\raggedright\arraybackslash}p{3.2cm}@{}}
\toprule
& \etiqueta{SQL Server} & \etiqueta{MySQL} & \etiqueta{SQLite} \\
\midrule
¿Hay servidor? & sí & sí & \textbf{no} \\
Programa para usarlo & SSMS & MySQL Workbench & DB Browser for SQLite \\
Su SQL se llama & Transact-SQL & SQL de MySQL & SQL de SQLite \\
Dónde se guarda & \texttt{.mdf} y \texttt{.ldf} & carpeta de datos; \texttt{.ibd} por tabla* & un archivo, p. ej.\ \texttt{.db} \\
Numerar solo el \texttt{id} & \texttt{IDENTITY(1,1)} & \texttt{AUTO\_INCREMENT} & \texttt{INTEGER PRIMARY KEY} \\
Las 5 primeras filas & \texttt{SELECT TOP 5 \ldots} & \texttt{\ldots{} LIMIT 5} & \texttt{\ldots{} LIMIT 5} \\
\bottomrule
\end{tabular}
\end{center}
\vspace{0.05cm}
{\scriptsize\color{utpGris} *MySQL: InnoDB con un archivo por tabla. SQLite puede crear archivos auxiliares.\\
Las 5 primeras filas requieren \texttt{ORDER BY} para definir un orden predecible.\par}
\vspace{0.15cm}
\small El modelo conceptual y el lógico de la peluquería \textbf{no cambian}. Para cada gestor solo se reescribe el \textbf{físico}.
\end{frame}
\begin{frame}[plain]
  \begin{tikzpicture}[remember picture,overlay]
    \node[anchor=north west,font=\bfseries\fontsize{24}{28}\selectfont]
      at ([xshift=12mm,yshift=-9mm]current page.north west) {Preguntas};
    \node[anchor=north east,inner sep=0pt]
      at ([xshift=-10mm,yshift=-8mm]current page.north east)
      {\includegraphics[width=43mm]{utp-logo.png}};
    \node[anchor=center,inner sep=0pt]
      at ([yshift=-5mm]current page.center)
      {\includegraphics[width=78mm]{preguntas-ilustracion.png}};
  \end{tikzpicture}
\end{frame}
\begin{frame}{Cierre: lo de hoy en una tabla}
\begin{center}
\small
\renewcommand{\arraystretch}{1.45}
\begin{tabular}{>{\raggedright\bfseries}p{4.2cm}>{\raggedright\arraybackslash}p{8.6cm}}
\toprule
& \bfseries En una frase \\
\midrule
Orden de \texttt{SELECT} & El orden escrito difiere del orden lógico. \\
Modelo conceptual & Entidades, atributos y relaciones del negocio. \\
Modelo lógico & Tablas, columnas, claves y restricciones. \\
Modelo físico & Tipos y SQL concretos para el gestor elegido. \\
Forward Engineering & Genera la estructura de la base a partir del modelo. \\
Reverse Engineering & Reconstruye un modelo a partir de una estructura existente. \\
Gestores & SQL Server y MySQL usan servidor. SQLite se integra en la aplicación. \\

\bottomrule
\end{tabular}
\end{center}
\vspace{0.25cm}
\small El conceptual y el lógico de la peluquería \textbf{se mantienen}. Adaptamos el \textbf{físico} a cada gestor.
\end{frame}

\begin{frame}{Referencias}
\footnotesize
\begin{itemize}\setlength\itemsep{10pt}
\item Coronel, C., Morris, S., \& Rob, P. (2011). \textit{Bases de datos: diseño, implementación y administración}. Cengage Learning.
\item Microsoft Learn. \href{https://learn.microsoft.com/en-us/sql/t-sql/queries/select-transact-sql}{\textit{SELECT (Transact-SQL): orden lógico de procesamiento}}.
\item Oracle. MySQL Workbench Manual: \href{https://dev.mysql.com/doc/workbench/en/wb-forward-engineering.html}{\textit{Forward Engineering}} y \href{https://dev.mysql.com/doc/workbench/en/wb-reverse-engineering.html}{\textit{Reverse Engineering}}.
\item Oracle. MySQL Reference Manual: \href{https://dev.mysql.com/doc/refman/8.4/en/innodb-file-per-table-tablespaces.html}{\textit{File-Per-Table Tablespaces}}.
\item Microsoft Learn. \href{https://learn.microsoft.com/en-us/sql/relational-databases/databases/database-files-and-filegroups}{\textit{Database Files and Filegroups}}.
\item SQLite. \href{https://www.sqlite.org/whentouse.html}{\textit{Appropriate Uses For SQLite}} y \href{https://www.sqlite.org/autoinc.html}{\textit{Autoincrement}}.
\end{itemize}
\end{frame}
% Cierre UTP
{
  \setbeamertemplate{background canvas}{%
    \includegraphics[width=\paperwidth,height=\paperheight]{gracias-final.pdf}%
  }
  \setbeamertemplate{footline}{}
  \begin{frame}[plain]
  \end{frame}
}

\end{document}
